\begin{abstract}
Failures of diabetic retinopathy (DR) grading models on new populations are usually blamed on the
learned representation, although ordinal models turn scores into grades with a rarely examined
decision rule. Under analysis plans committed in advance, we evaluated ten development models on four external populations (APTOS 2019, DDR,
DeepDRiD, SUSTech-SYSU), a reverse-direction model, two models fine-tuned
on \pn{eye.eb3.trainn} EyePACS images, five seeds and a published model's released predictions,
separating per-threshold discrimination from the score-to-grade mapping.
Referral ranking largely transferred (area under the curve at least \pn{r2.auc1.allmin}), but the
score given the true grade moved in opposite directions: the same models over-graded APTOS and
under-graded DDR, so a threshold fitted for 85\% development sensitivity reached
\pn{aptos.tr.e08.B.sens}\% and \pn{ddr.tr.e08.B.sens}\%. Fitted on EyePACS instead, the threshold
over-referred everywhere (\pn{eye.tr85.min}--\pn{eye.tr85.max}\%): the direction of the shift
belonged to the population, the sign of the threshold error to the development labels. Refitting
cut-points recovered mild-grade recall on APTOS and reverse-direction referral sensitivity
(\pn{rev.m2.R1.sens}\% to \pn{rev.m2.R4.sens}\%) but not lost ranking; label-free prior adaptation
did not help. Contrary to two predictions, models differed more than deployable rules, and the
collapse of mild recall was specific to one corpus, although mild F1 stayed low everywhere. The
risk that a locally set threshold misses an 80\% sensitivity floor depended on the number of
referable cases labelled, not of images: about 120 kept it near one in ten (about 1\,000 images at
10\% prevalence). External evaluations should report per-threshold discrimination and operating
curves, and referral thresholds should be set locally.
\end{abstract}
